\documentclass[acmsmall,screen,review,anonymous,nonacm]{acmart}

\usepackage{paper}

\begin{document}

% Generated by experiments/paper_numbers.py from results/ -- do not edit.
\newcommand{\NBareboxInstances}{191\xspace}
\newcommand{\NBareboxMembers}{89\xspace}
\newcommand{\NBareboxNoIncludesLostObjects}{0\xspace}
\newcommand{\NBareboxNoIncludesSandboxDefconfigFn}{0\xspace}
\newcommand{\NBareboxNoIncludesSandboxDefconfigFp}{0\xspace}
\newcommand{\NBareboxNoIncludesSandboxDefconfigOverlap}{92.0\xspace}
\newcommand{\NBareboxNoLinkSemanticsLostObjects}{-28\xspace}
\newcommand{\NBareboxNoLinkSemanticsSandboxDefconfigFn}{0\xspace}
\newcommand{\NBareboxNoLinkSemanticsSandboxDefconfigFp}{4\xspace}
\newcommand{\NBareboxNoLinkSemanticsSandboxDefconfigOverlap}{92.0\xspace}
\newcommand{\NBareboxNoRulesLostObjects}{0\xspace}
\newcommand{\NBareboxNoRulesSandboxDefconfigFn}{0\xspace}
\newcommand{\NBareboxNoRulesSandboxDefconfigFp}{0\xspace}
\newcommand{\NBareboxNoRulesSandboxDefconfigOverlap}{92.0\xspace}
\newcommand{\NBareboxObjects}{1,770\xspace}
\newcommand{\NBareboxRss}{70.3\xspace}
\newcommand{\NBareboxSandboxDefconfigBuildMinutes}{0.3\xspace}
\newcommand{\NBareboxSandboxDefconfigFn}{0\xspace}
\newcommand{\NBareboxSandboxDefconfigFp}{0\xspace}
\newcommand{\NBareboxSandboxDefconfigFpbuild}{0\xspace}
\newcommand{\NBareboxSandboxDefconfigOut}{40\xspace}
\newcommand{\NBareboxSandboxDefconfigOuthost}{32\xspace}
\newcommand{\NBareboxSandboxDefconfigOuttarget}{8\xspace}
\newcommand{\NBareboxSandboxDefconfigOverlap}{92.0\xspace}
\newcommand{\NBareboxSandboxDefconfigOverlapnohost}{98.3\xspace}
\newcommand{\NBareboxSandboxDefconfigPhysical}{502\xspace}
\newcommand{\NBareboxSandboxDefconfigPhysicalObjects}{502\xspace}
\newcommand{\NBareboxSandboxDefconfigPrec}{100.0\xspace}
\newcommand{\NBareboxSandboxDefconfigPred}{462\xspace}
\newcommand{\NBareboxSandboxDefconfigTargetsoverlap}{89.8\xspace}
\newcommand{\NBareboxSandboxDefconfigTp}{462\xspace}
\newcommand{\NBareboxSeconds}{2.7\xspace}
\newcommand{\NBareboxTargets}{1,681\xspace}
\newcommand{\NBareboxVersion}{2026.09.0\xspace}
\newcommand{\NBareboxZTchecks}{4,948\xspace}
\newcommand{\NBenchAppendESeconds}{0.5\xspace}
\newcommand{\NBenchAppendEStates}{16\xspace}
\newcommand{\NBenchAppendFESeconds}{0.5\xspace}
\newcommand{\NBenchAppendFEStates}{96\xspace}
\newcommand{\NBenchAppendFSeconds}{0.5\xspace}
\newcommand{\NBenchAppendFStates}{8\xspace}
\newcommand{\NBenchAppendOSSeconds}{0.5\xspace}
\newcommand{\NBenchAppendOSStates}{32\xspace}
\newcommand{\NBenchAppendOSeconds}{0.6\xspace}
\newcommand{\NBenchAppendOStates}{2\xspace}
\newcommand{\NBenchAppendOTESeconds}{0.6\xspace}
\newcommand{\NBenchAppendOTEStates}{256\xspace}
\newcommand{\NBenchAppendOTSeconds}{0.5\xspace}
\newcommand{\NBenchAppendOTStates}{24\xspace}
\newcommand{\NBenchAppendSFSeconds}{0.5\xspace}
\newcommand{\NBenchAppendSFStates}{128\xspace}
\newcommand{\NBenchAppendTFSeconds}{0.6\xspace}
\newcommand{\NBenchAppendTFStates}{48\xspace}
\newcommand{\NBenchAppendTSeconds}{0.5\xspace}
\newcommand{\NBenchAppendTStates}{4\xspace}
\newcommand{\NBenchAppendTTSeconds}{0.5\xspace}
\newcommand{\NBenchAppendTTStates}{64\xspace}
\newcommand{\NBenchAppendTZSeconds}{0.5\xspace}
\newcommand{\NBenchAppendTZStates}{40\xspace}
\newcommand{\NBenchIfeqESeconds}{0.7\xspace}
\newcommand{\NBenchIfeqEStates}{9\xspace}
\newcommand{\NBenchIfeqFESeconds}{68.3\xspace}
\newcommand{\NBenchIfeqFEStates}{49\xspace}
\newcommand{\NBenchIfeqFSeconds}{0.6\xspace}
\newcommand{\NBenchIfeqFStates}{5\xspace}
\newcommand{\NBenchIfeqOSSeconds}{1.5\xspace}
\newcommand{\NBenchIfeqOSStates}{17\xspace}
\newcommand{\NBenchIfeqOSeconds}{0.5\xspace}
\newcommand{\NBenchIfeqOStates}{2\xspace}
\newcommand{\NBenchIfeqOTSeconds}{1.0\xspace}
\newcommand{\NBenchIfeqOTStates}{13\xspace}
\newcommand{\NBenchIfeqSFSeconds}{268.4\xspace}
\newcommand{\NBenchIfeqSFStates}{65\xspace}
\newcommand{\NBenchIfeqTFSeconds}{4.3\xspace}
\newcommand{\NBenchIfeqTFStates}{25\xspace}
\newcommand{\NBenchIfeqTSeconds}{0.6\xspace}
\newcommand{\NBenchIfeqTStates}{3\xspace}
\newcommand{\NBenchIfeqTTSeconds}{12.3\xspace}
\newcommand{\NBenchIfeqTTStates}{33\xspace}
\newcommand{\NBenchIfeqTZSeconds}{2.5\xspace}
\newcommand{\NBenchIfeqTZStates}{21\xspace}
\newcommand{\NBenchIndirectESeconds}{0.9\xspace}
\newcommand{\NBenchIndirectEStates}{1,025\xspace}
\newcommand{\NBenchIndirectFSeconds}{0.5\xspace}
\newcommand{\NBenchIndirectFStates}{33\xspace}
\newcommand{\NBenchIndirectOSSeconds}{187.1\xspace}
\newcommand{\NBenchIndirectOSStates}{524,289\xspace}
\newcommand{\NBenchIndirectOSeconds}{0.5\xspace}
\newcommand{\NBenchIndirectOStates}{2\xspace}
\newcommand{\NBenchIndirectOTSeconds}{8.7\xspace}
\newcommand{\NBenchIndirectOTStates}{24,577\xspace}
\newcommand{\NBenchIndirectTSeconds}{0.5\xspace}
\newcommand{\NBenchIndirectTStates}{5\xspace}
\newcommand{\NBenchKfoldInexact}{0\xspace}
\newcommand{\NBenchKfoldMaxRss}{67\xspace}
\newcommand{\NBenchKfoldMaxSeconds}{0.73\xspace}
\newcommand{\NBenchKfoldMinSeconds}{0.61\xspace}
\newcommand{\NBenchKfoldRuns}{48\xspace}
\newcommand{\NBenchKmaxInexact}{12\xspace}
\newcommand{\NBenchKmaxMaxRss}{73\xspace}
\newcommand{\NBenchKmaxMaxSeconds}{0.92\xspace}
\newcommand{\NBenchKmaxMinSeconds}{0.67\xspace}
\newcommand{\NBenchKmaxRuns}{48\xspace}
\newcommand{\NBenchKmaxSmtInexact}{12\xspace}
\newcommand{\NBenchKmaxSmtMaxRss}{134\xspace}
\newcommand{\NBenchKmaxSmtMaxSeconds}{6.78\xspace}
\newcommand{\NBenchKmaxSmtMinSeconds}{0.68\xspace}
\newcommand{\NBenchKmaxSmtRuns}{48\xspace}
\newcommand{\NBenchMaxk}{128\xspace}
\newcommand{\NBenchTimeout}{300\xspace}
\newcommand{\NBenchValueESeconds}{3.2\xspace}
\newcommand{\NBenchValueEStates}{2,048\xspace}
\newcommand{\NBenchValueFSeconds}{0.6\xspace}
\newcommand{\NBenchValueFStates}{64\xspace}
\newcommand{\NBenchValueOSeconds}{0.5\xspace}
\newcommand{\NBenchValueOStates}{2\xspace}
\newcommand{\NBenchValueOTSeconds}{77.8\xspace}
\newcommand{\NBenchValueOTStates}{49,152\xspace}
\newcommand{\NBenchValueTSeconds}{0.5\xspace}
\newcommand{\NBenchValueTStates}{8\xspace}
\newcommand{\NBlindArch}{349\xspace}
\newcommand{\NBlindBlindSpots}{3,199\xspace}
\newcommand{\NBlindEntriesWithBlindSpots}{509\xspace}
\newcommand{\NBlindNopivot}{664\xspace}
\newcommand{\NBlindNotenabled}{2,186\xspace}
\newcommand{\NBlindObjectsConsidered}{32,829\xspace}
\newcommand{\NBusyboxDefconfigBuildMinutes}{0.3\xspace}
\newcommand{\NBusyboxDefconfigFn}{0\xspace}
\newcommand{\NBusyboxDefconfigFp}{0\xspace}
\newcommand{\NBusyboxDefconfigFpbuild}{1\xspace}
\newcommand{\NBusyboxDefconfigOut}{34\xspace}
\newcommand{\NBusyboxDefconfigOuthost}{4\xspace}
\newcommand{\NBusyboxDefconfigOuttarget}{30\xspace}
\newcommand{\NBusyboxDefconfigOverlap}{94.5\xspace}
\newcommand{\NBusyboxDefconfigOverlapnohost}{95.2\xspace}
\newcommand{\NBusyboxDefconfigPhysical}{623\xspace}
\newcommand{\NBusyboxDefconfigPhysicalObjects}{623\xspace}
\newcommand{\NBusyboxDefconfigPrec}{100.0\xspace}
\newcommand{\NBusyboxDefconfigPred}{590\xspace}
\newcommand{\NBusyboxDefconfigTargetsoverlap}{94.5\xspace}
\newcommand{\NBusyboxDefconfigTp}{589\xspace}
\newcommand{\NBusyboxInstances}{30\xspace}
\newcommand{\NBusyboxKmaxDefconfigLocalFn}{0\xspace}
\newcommand{\NBusyboxKmaxDefconfigLocalOverlap}{94.5\xspace}
\newcommand{\NBusyboxKmaxDefconfigLocalPrec}{99.8\xspace}
\newcommand{\NBusyboxKmaxDefconfigLocalTp}{589\xspace}
\newcommand{\NBusyboxKmaxDefconfigWithDirsFn}{0\xspace}
\newcommand{\NBusyboxKmaxDefconfigWithDirsOverlap}{94.5\xspace}
\newcommand{\NBusyboxKmaxDefconfigWithDirsPrec}{99.8\xspace}
\newcommand{\NBusyboxKmaxDefconfigWithDirsTp}{589\xspace}
\newcommand{\NBusyboxKmaxObjects}{625\xspace}
\newcommand{\NBusyboxKmaxRssGb}{0.1\xspace}
\newcommand{\NBusyboxKmaxSeconds}{73\xspace}
\newcommand{\NBusyboxKmaxUnexpanded}{0\xspace}
\newcommand{\NBusyboxMembers}{0\xspace}
\newcommand{\NBusyboxNoIncludesDefconfigFn}{0\xspace}
\newcommand{\NBusyboxNoIncludesDefconfigFp}{0\xspace}
\newcommand{\NBusyboxNoIncludesDefconfigOverlap}{94.5\xspace}
\newcommand{\NBusyboxNoIncludesLostObjects}{0\xspace}
\newcommand{\NBusyboxNoLinkSemanticsDefconfigFn}{0\xspace}
\newcommand{\NBusyboxNoLinkSemanticsDefconfigFp}{0\xspace}
\newcommand{\NBusyboxNoLinkSemanticsDefconfigOverlap}{94.5\xspace}
\newcommand{\NBusyboxNoLinkSemanticsLostObjects}{0\xspace}
\newcommand{\NBusyboxNoRulesDefconfigFn}{0\xspace}
\newcommand{\NBusyboxNoRulesDefconfigFp}{0\xspace}
\newcommand{\NBusyboxNoRulesDefconfigOverlap}{94.5\xspace}
\newcommand{\NBusyboxNoRulesLostObjects}{0\xspace}
\newcommand{\NBusyboxObjects}{625\xspace}
\newcommand{\NBusyboxRss}{68.8\xspace}
\newcommand{\NBusyboxSeconds}{0.3\xspace}
\newcommand{\NBusyboxTargets}{625\xspace}
\newcommand{\NBusyboxVersion}{1.38.0\xspace}
\newcommand{\NBusyboxZTchecks}{2\xspace}
\newcommand{\NCompileOk}{40\xspace}
\newcommand{\NCompileSampled}{40\xspace}
\newcommand{\NConfigforCommits}{402\xspace}
\newcommand{\NConfigforExcluded}{6\xspace}
\newcommand{\NConfigforExcludedObjects}{6\xspace}
\newcommand{\NConfigforLost}{0\xspace}
\newcommand{\NConfigforMedianSeconds}{6.7\xspace}
\newcommand{\NConfigforObjects}{471\xspace}
\newcommand{\NConfigforSelectable}{384\xspace}
\newcommand{\NConfigforUnmapped}{12\xspace}
\newcommand{\NCorebootInstances}{513\xspace}
\newcommand{\NCorebootMembers}{0\xspace}
\newcommand{\NCorebootNoIncludesLostObjects}{0\xspace}
\newcommand{\NCorebootNoIncludesQemuIFFZfxFn}{0\xspace}
\newcommand{\NCorebootNoIncludesQemuIFFZfxFp}{0\xspace}
\newcommand{\NCorebootNoIncludesQemuIFFZfxOverlap}{76.2\xspace}
\newcommand{\NCorebootNoLinkSemanticsLostObjects}{0\xspace}
\newcommand{\NCorebootNoLinkSemanticsQemuIFFZfxFn}{0\xspace}
\newcommand{\NCorebootNoLinkSemanticsQemuIFFZfxFp}{0\xspace}
\newcommand{\NCorebootNoLinkSemanticsQemuIFFZfxOverlap}{76.2\xspace}
\newcommand{\NCorebootNoRulesLostObjects}{0\xspace}
\newcommand{\NCorebootNoRulesQemuIFFZfxFn}{0\xspace}
\newcommand{\NCorebootNoRulesQemuIFFZfxFp}{0\xspace}
\newcommand{\NCorebootNoRulesQemuIFFZfxOverlap}{76.2\xspace}
\newcommand{\NCorebootObjects}{5,444\xspace}
\newcommand{\NCorebootQemuIFFZfxBuildMinutes}{0.9\xspace}
\newcommand{\NCorebootQemuIFFZfxFn}{0\xspace}
\newcommand{\NCorebootQemuIFFZfxFp}{0\xspace}
\newcommand{\NCorebootQemuIFFZfxFpbuild}{0\xspace}
\newcommand{\NCorebootQemuIFFZfxOut}{149\xspace}
\newcommand{\NCorebootQemuIFFZfxOuthost}{140\xspace}
\newcommand{\NCorebootQemuIFFZfxOuttarget}{9\xspace}
\newcommand{\NCorebootQemuIFFZfxOverlap}{76.2\xspace}
\newcommand{\NCorebootQemuIFFZfxOverlapnohost}{98.1\xspace}
\newcommand{\NCorebootQemuIFFZfxPhysical}{626\xspace}
\newcommand{\NCorebootQemuIFFZfxPhysicalObjects}{626\xspace}
\newcommand{\NCorebootQemuIFFZfxPrec}{100.0\xspace}
\newcommand{\NCorebootQemuIFFZfxPred}{477\xspace}
\newcommand{\NCorebootQemuIFFZfxTargetsoverlap}{76.2\xspace}
\newcommand{\NCorebootQemuIFFZfxTp}{477\xspace}
\newcommand{\NCorebootRss}{106.1\xspace}
\newcommand{\NCorebootSeconds}{23.0\xspace}
\newcommand{\NCorebootTargets}{5,444\xspace}
\newcommand{\NCorebootVersion}{26.06\xspace}
\newcommand{\NCorebootZTchecks}{4,163\xspace}
\newcommand{\NKconfigAllmodconfigMismatches}{1\xspace}
\newcommand{\NKconfigAllmodconfigSymbols}{16,102\xspace}
\newcommand{\NKconfigAllyesconfigMismatches}{0\xspace}
\newcommand{\NKconfigAllyesconfigSymbols}{16,208\xspace}
\newcommand{\NKconfigConfigs}{5\xspace}
\newcommand{\NKconfigDebianMismatches}{0\xspace}
\newcommand{\NKconfigDebianSymbols}{9,881\xspace}
\newcommand{\NKconfigDefconfigMismatches}{0\xspace}
\newcommand{\NKconfigDefconfigSymbols}{4,414\xspace}
\newcommand{\NKconfigDefined}{18,749\xspace}
\newcommand{\NKconfigMismatches}{1\xspace}
\newcommand{\NKconfigSmtViolations}{0\xspace}
\newcommand{\NKconfigTinyconfigMismatches}{0\xspace}
\newcommand{\NKconfigTinyconfigSymbols}{859\xspace}
\newcommand{\NLinuxAllmodconfigBuildMinutes}{100.2\xspace}
\newcommand{\NLinuxAllmodconfigFn}{0\xspace}
\newcommand{\NLinuxAllmodconfigFp}{0\xspace}
\newcommand{\NLinuxAllmodconfigFpbuild}{0\xspace}
\newcommand{\NLinuxAllmodconfigOut}{69\xspace}
\newcommand{\NLinuxAllmodconfigOuthost}{69\xspace}
\newcommand{\NLinuxAllmodconfigOuttarget}{0\xspace}
\newcommand{\NLinuxAllmodconfigOverlap}{99.8\xspace}
\newcommand{\NLinuxAllmodconfigOverlapnohost}{100.0\xspace}
\newcommand{\NLinuxAllmodconfigPhysical}{29,625\xspace}
\newcommand{\NLinuxAllmodconfigPhysicalObjects}{29,625\xspace}
\newcommand{\NLinuxAllmodconfigPrec}{100.0\xspace}
\newcommand{\NLinuxAllmodconfigPred}{29,556\xspace}
\newcommand{\NLinuxAllmodconfigTargetsoverlap}{47.0\xspace}
\newcommand{\NLinuxAllmodconfigTp}{29,556\xspace}
\newcommand{\NLinuxDebianBuildMinutes}{54.2\xspace}
\newcommand{\NLinuxDebianFn}{0\xspace}
\newcommand{\NLinuxDebianFp}{0\xspace}
\newcommand{\NLinuxDebianFpbuild}{0\xspace}
\newcommand{\NLinuxDebianOut}{91\xspace}
\newcommand{\NLinuxDebianOuthost}{91\xspace}
\newcommand{\NLinuxDebianOuttarget}{0\xspace}
\newcommand{\NLinuxDebianOverlap}{99.5\xspace}
\newcommand{\NLinuxDebianOverlapnohost}{100.0\xspace}
\newcommand{\NLinuxDebianPhysical}{17,191\xspace}
\newcommand{\NLinuxDebianPhysicalObjects}{17,191\xspace}
\newcommand{\NLinuxDebianPrec}{100.0\xspace}
\newcommand{\NLinuxDebianPred}{17,100\xspace}
\newcommand{\NLinuxDebianTargetsoverlap}{35.7\xspace}
\newcommand{\NLinuxDebianTp}{17,100\xspace}
\newcommand{\NLinuxDefconfigBuildMinutes}{16.9\xspace}
\newcommand{\NLinuxDefconfigFn}{0\xspace}
\newcommand{\NLinuxDefconfigFp}{0\xspace}
\newcommand{\NLinuxDefconfigFpbuild}{0\xspace}
\newcommand{\NLinuxDefconfigOut}{48\xspace}
\newcommand{\NLinuxDefconfigOuthost}{48\xspace}
\newcommand{\NLinuxDefconfigOuttarget}{0\xspace}
\newcommand{\NLinuxDefconfigOverlap}{98.5\xspace}
\newcommand{\NLinuxDefconfigOverlapnohost}{100.0\xspace}
\newcommand{\NLinuxDefconfigPhysical}{3,133\xspace}
\newcommand{\NLinuxDefconfigPhysicalObjects}{3,133\xspace}
\newcommand{\NLinuxDefconfigPrec}{100.0\xspace}
\newcommand{\NLinuxDefconfigPred}{3,085\xspace}
\newcommand{\NLinuxDefconfigTargetsoverlap}{47.2\xspace}
\newcommand{\NLinuxDefconfigTp}{3,085\xspace}
\newcommand{\NLinuxInstances}{2,215\xspace}
\newcommand{\NLinuxKmaxAllmodconfigLocalFn}{8,600\xspace}
\newcommand{\NLinuxKmaxAllmodconfigLocalFnMembers}{8,549\xspace}
\newcommand{\NLinuxKmaxAllmodconfigLocalOverlap}{63.0\xspace}
\newcommand{\NLinuxKmaxAllmodconfigLocalPrec}{93.7\xspace}
\newcommand{\NLinuxKmaxAllmodconfigLocalTp}{18,656\xspace}
\newcommand{\NLinuxKmaxAllmodconfigWithDirsFn}{9,038\xspace}
\newcommand{\NLinuxKmaxAllmodconfigWithDirsFnMembers}{8,576\xspace}
\newcommand{\NLinuxKmaxAllmodconfigWithDirsOverlap}{61.5\xspace}
\newcommand{\NLinuxKmaxAllmodconfigWithDirsPrec}{94.6\xspace}
\newcommand{\NLinuxKmaxAllmodconfigWithDirsTp}{18,218\xspace}
\newcommand{\NLinuxKmaxDebianLocalFn}{6,138\xspace}
\newcommand{\NLinuxKmaxDebianLocalFnMembers}{6,087\xspace}
\newcommand{\NLinuxKmaxDebianLocalOverlap}{56.3\xspace}
\newcommand{\NLinuxKmaxDebianLocalPrec}{88.6\xspace}
\newcommand{\NLinuxKmaxDebianLocalTp}{9,675\xspace}
\newcommand{\NLinuxKmaxDebianWithDirsFn}{6,471\xspace}
\newcommand{\NLinuxKmaxDebianWithDirsFnMembers}{6,108\xspace}
\newcommand{\NLinuxKmaxDebianWithDirsOverlap}{54.3\xspace}
\newcommand{\NLinuxKmaxDebianWithDirsPrec}{90.1\xspace}
\newcommand{\NLinuxKmaxDebianWithDirsTp}{9,342\xspace}
\newcommand{\NLinuxKmaxDefconfigLocalFn}{634\xspace}
\newcommand{\NLinuxKmaxDefconfigLocalFnMembers}{597\xspace}
\newcommand{\NLinuxKmaxDefconfigLocalOverlap}{73.0\xspace}
\newcommand{\NLinuxKmaxDefconfigLocalPrec}{77.2\xspace}
\newcommand{\NLinuxKmaxDefconfigLocalTp}{2,286\xspace}
\newcommand{\NLinuxKmaxDefconfigWithDirsFn}{834\xspace}
\newcommand{\NLinuxKmaxDefconfigWithDirsFnMembers}{598\xspace}
\newcommand{\NLinuxKmaxDefconfigWithDirsOverlap}{66.6\xspace}
\newcommand{\NLinuxKmaxDefconfigWithDirsPrec}{85.5\xspace}
\newcommand{\NLinuxKmaxDefconfigWithDirsTp}{2,086\xspace}
\newcommand{\NLinuxKmaxObjects}{31,751\xspace}
\newcommand{\NLinuxKmaxRssGb}{27.9\xspace}
\newcommand{\NLinuxKmaxSeconds}{6,275\xspace}
\newcommand{\NLinuxKmaxTinyconfigLocalFn}{27\xspace}
\newcommand{\NLinuxKmaxTinyconfigLocalFnMembers}{5\xspace}
\newcommand{\NLinuxKmaxTinyconfigLocalOverlap}{78.0\xspace}
\newcommand{\NLinuxKmaxTinyconfigLocalPrec}{31.4\xspace}
\newcommand{\NLinuxKmaxTinyconfigLocalTp}{412\xspace}
\newcommand{\NLinuxKmaxTinyconfigWithDirsFn}{129\xspace}
\newcommand{\NLinuxKmaxTinyconfigWithDirsFnMembers}{5\xspace}
\newcommand{\NLinuxKmaxTinyconfigWithDirsOverlap}{58.7\xspace}
\newcommand{\NLinuxKmaxTinyconfigWithDirsPrec}{96.3\xspace}
\newcommand{\NLinuxKmaxTinyconfigWithDirsTp}{310\xspace}
\newcommand{\NLinuxKmaxUnexpanded}{28\xspace}
\newcommand{\NLinuxMembers}{16,336\xspace}
\newcommand{\NLinuxMinOverlapnohost}{100.0\xspace}
\newcommand{\NLinuxNoIncludesAllmodconfigFn}{0\xspace}
\newcommand{\NLinuxNoIncludesAllmodconfigFp}{0\xspace}
\newcommand{\NLinuxNoIncludesAllmodconfigLostAmd}{687\xspace}
\newcommand{\NLinuxNoIncludesAllmodconfigLostNouveau}{732\xspace}
\newcommand{\NLinuxNoIncludesAllmodconfigOverlap}{94.7\xspace}
\newcommand{\NLinuxNoIncludesDebianFn}{0\xspace}
\newcommand{\NLinuxNoIncludesDebianFp}{0\xspace}
\newcommand{\NLinuxNoIncludesDebianOverlap}{90.8\xspace}
\newcommand{\NLinuxNoIncludesDefconfigFn}{0\xspace}
\newcommand{\NLinuxNoIncludesDefconfigFp}{0\xspace}
\newcommand{\NLinuxNoIncludesDefconfigOverlap}{98.0\xspace}
\newcommand{\NLinuxNoIncludesLostObjects}{1,507\xspace}
\newcommand{\NLinuxNoIncludesTinyconfigFn}{0\xspace}
\newcommand{\NLinuxNoIncludesTinyconfigFp}{0\xspace}
\newcommand{\NLinuxNoIncludesTinyconfigOverlap}{95.8\xspace}
\newcommand{\NLinuxNoLinkSemanticsAllmodconfigFn}{0\xspace}
\newcommand{\NLinuxNoLinkSemanticsAllmodconfigFp}{144\xspace}
\newcommand{\NLinuxNoLinkSemanticsAllmodconfigOverlap}{99.8\xspace}
\newcommand{\NLinuxNoLinkSemanticsDebianFn}{0\xspace}
\newcommand{\NLinuxNoLinkSemanticsDebianFp}{121\xspace}
\newcommand{\NLinuxNoLinkSemanticsDebianOverlap}{99.5\xspace}
\newcommand{\NLinuxNoLinkSemanticsDefconfigFn}{0\xspace}
\newcommand{\NLinuxNoLinkSemanticsDefconfigFp}{161\xspace}
\newcommand{\NLinuxNoLinkSemanticsDefconfigOverlap}{98.5\xspace}
\newcommand{\NLinuxNoLinkSemanticsLostObjects}{-14\xspace}
\newcommand{\NLinuxNoLinkSemanticsTinyconfigFn}{0\xspace}
\newcommand{\NLinuxNoLinkSemanticsTinyconfigFp}{5\xspace}
\newcommand{\NLinuxNoLinkSemanticsTinyconfigOverlap}{97.0\xspace}
\newcommand{\NLinuxNoRulesAllmodconfigFn}{10\xspace}
\newcommand{\NLinuxNoRulesAllmodconfigFp}{0\xspace}
\newcommand{\NLinuxNoRulesAllmodconfigOverlap}{99.4\xspace}
\newcommand{\NLinuxNoRulesDebianFn}{10\xspace}
\newcommand{\NLinuxNoRulesDebianFp}{0\xspace}
\newcommand{\NLinuxNoRulesDebianOverlap}{98.8\xspace}
\newcommand{\NLinuxNoRulesDefconfigFn}{5\xspace}
\newcommand{\NLinuxNoRulesDefconfigFp}{0\xspace}
\newcommand{\NLinuxNoRulesDefconfigOverlap}{95.8\xspace}
\newcommand{\NLinuxNoRulesLostObjects}{112\xspace}
\newcommand{\NLinuxNoRulesTinyconfigFn}{0\xspace}
\newcommand{\NLinuxNoRulesTinyconfigFp}{0\xspace}
\newcommand{\NLinuxNoRulesTinyconfigOverlap}{89.2\xspace}
\newcommand{\NLinuxObjects}{32,829\xspace}
\newcommand{\NLinuxRss}{203.9\xspace}
\newcommand{\NLinuxSeconds}{143.1\xspace}
\newcommand{\NLinuxTargets}{16,366\xspace}
\newcommand{\NLinuxTinyconfigBuildMinutes}{2.6\xspace}
\newcommand{\NLinuxTinyconfigFn}{0\xspace}
\newcommand{\NLinuxTinyconfigFp}{0\xspace}
\newcommand{\NLinuxTinyconfigFpbuild}{0\xspace}
\newcommand{\NLinuxTinyconfigOut}{16\xspace}
\newcommand{\NLinuxTinyconfigOuthost}{16\xspace}
\newcommand{\NLinuxTinyconfigOuttarget}{0\xspace}
\newcommand{\NLinuxTinyconfigOverlap}{97.0\xspace}
\newcommand{\NLinuxTinyconfigOverlapnohost}{100.0\xspace}
\newcommand{\NLinuxTinyconfigPhysical}{528\xspace}
\newcommand{\NLinuxTinyconfigPhysicalObjects}{528\xspace}
\newcommand{\NLinuxTinyconfigPrec}{100.0\xspace}
\newcommand{\NLinuxTinyconfigPred}{512\xspace}
\newcommand{\NLinuxTinyconfigTargetsoverlap}{86.6\xspace}
\newcommand{\NLinuxTinyconfigTp}{512\xspace}
\newcommand{\NLinuxVersion}{7.2.8\xspace}
\newcommand{\NLinuxZTchecks}{107,208\xspace}
\newcommand{\NUbootInstances}{319\xspace}
\newcommand{\NUbootMembers}{128\xspace}
\newcommand{\NUbootNoIncludesLostObjects}{0\xspace}
\newcommand{\NUbootNoIncludesSandboxDefconfigFn}{0\xspace}
\newcommand{\NUbootNoIncludesSandboxDefconfigFp}{0\xspace}
\newcommand{\NUbootNoIncludesSandboxDefconfigOverlap}{90.3\xspace}
\newcommand{\NUbootNoLinkSemanticsLostObjects}{-21\xspace}
\newcommand{\NUbootNoLinkSemanticsSandboxDefconfigFn}{0\xspace}
\newcommand{\NUbootNoLinkSemanticsSandboxDefconfigFp}{9\xspace}
\newcommand{\NUbootNoLinkSemanticsSandboxDefconfigOverlap}{90.3\xspace}
\newcommand{\NUbootNoRulesLostObjects}{9\xspace}
\newcommand{\NUbootNoRulesSandboxDefconfigFn}{0\xspace}
\newcommand{\NUbootNoRulesSandboxDefconfigFp}{0\xspace}
\newcommand{\NUbootNoRulesSandboxDefconfigOverlap}{90.3\xspace}
\newcommand{\NUbootObjects}{3,750\xspace}
\newcommand{\NUbootRss}{77.6\xspace}
\newcommand{\NUbootSandboxDefconfigBuildMinutes}{0.7\xspace}
\newcommand{\NUbootSandboxDefconfigFn}{0\xspace}
\newcommand{\NUbootSandboxDefconfigFp}{0\xspace}
\newcommand{\NUbootSandboxDefconfigFpbuild}{0\xspace}
\newcommand{\NUbootSandboxDefconfigOut}{130\xspace}
\newcommand{\NUbootSandboxDefconfigOuthost}{115\xspace}
\newcommand{\NUbootSandboxDefconfigOuttarget}{15\xspace}
\newcommand{\NUbootSandboxDefconfigOverlap}{90.3\xspace}
\newcommand{\NUbootSandboxDefconfigOverlapnohost}{98.8\xspace}
\newcommand{\NUbootSandboxDefconfigPhysical}{1,344\xspace}
\newcommand{\NUbootSandboxDefconfigPhysicalObjects}{1,344\xspace}
\newcommand{\NUbootSandboxDefconfigPrec}{100.0\xspace}
\newcommand{\NUbootSandboxDefconfigPred}{1,214\xspace}
\newcommand{\NUbootSandboxDefconfigTargetsoverlap}{86.9\xspace}
\newcommand{\NUbootSandboxDefconfigTp}{1,214\xspace}
\newcommand{\NUbootSeconds}{4.8\xspace}
\newcommand{\NUbootTargets}{3,611\xspace}
\newcommand{\NUbootVersion}{2026.07\xspace}
\newcommand{\NUbootZTchecks}{7,895\xspace}
\newcommand{\NWhyCases}{400\xspace}
\newcommand{\NWhyConfigs}{4\xspace}
\newcommand{\NWhyCorrect}{400\xspace}
\newcommand{\NWhyFailCompositeMembership}{20\xspace}
\newcommand{\NWhyFailDirectoryRoute}{184\xspace}
\newcommand{\NWhyFailTargetListLine}{63\xspace}
\newcommand{\NWhyPerSide}{50\xspace}


\title{Variability-Aware Analysis of Kbuild Makefiles}

\author{Anonymous Author(s)}
\affiliation{\institution{Anonymous Institution} \country{}}
\renewcommand{\shortauthors}{Anonymous Author(s)}

\begin{abstract}
  The Linux kernel is highly configurable, and its \texttt{kbuild} Makefiles decide how thousands of configuration options affect which source files are included in the kernel build.
  A developer who asks why a source file is not built, or which options would build it, needs its build condition, a logical formula over configuration options.
  Running Make reveals the files of one configuration only, and enumerating all configurations is intractable.
  Symbolic execution that forks a state at every conditional suffers from path explosion, and prior static analyses of Kbuild Makefiles represent a variable as a set of complete string values.
  We present \tool{}, a symbolic execution of \texttt{kbuild} Makefiles that maps each Make variable to a set of words, each guarded by a path condition over configuration options, in the style of a value summary.
  Appends add guarded words, an overwrite re-guards the words it replaces, and path merging joins the two arms of a conditional into one symbolic state, so the state grows with the number of words instead of the number of paths.
  The resulting build conditions are SMT formulae that a solver can check and satisfy.

  We evaluate \tool{} on Linux \NLinuxVersion{} and four other \texttt{kbuild}-based systems by comparing its build conditions under concrete configurations with the source files built by fresh builds.
  On four Linux configurations, from \texttt{tinyconfig} to \texttt{allmodconfig}, the predicted files agree with the built files, with no false positive and no miss, from \NLinuxTinyconfigTp{} to \NLinuxAllmodconfigTp{} files.
  \tool{} analyzes the \NLinuxInstances{} Makefile instances of the Linux x86 tree in \NLinuxSeconds{}\,s.
  The build conditions explain the build status of all \NWhyCases{} sampled Linux files and, combined with a Kconfig model, produce configurations for \NConfigforSelectable{} of the \NConfigforCommits{} commits of a Linux stable update that touch C or assembly files.
\end{abstract}

\begin{CCSXML}
<ccs2012>
<concept>
<concept_id>10011007.10011074.10011099</concept_id>
<concept_desc>Software and its engineering~Software verification and validation</concept_desc>
<concept_significance>500</concept_significance>
</concept>
</ccs2012>
\end{CCSXML}
\ccsdesc[500]{Software and its engineering~Software verification and validation}
\keywords{Kbuild, Makefile, symbolic execution, variability, Kconfig}

\maketitle

\section{Introduction}\label{sec:intro}
\placeholder{} \cite{placeholder2027}

\section{Background}\label{sec:background}

\section{Motivating Example}\label{sec:example}

\section{Approach}\label{sec:approach}

\section{Implementation}\label{sec:impl}

\section{Evaluation}\label{sec:eval}
\subsection{Experimental Setup}\label{sec:setup}
\subsection{RQ1}\label{sec:rq1}
\subsection{RQ2}\label{sec:rq2}

\section{Related Work}\label{sec:related}

\section{Conclusion}\label{sec:conclusion}

\section*{Data-Availability Statement}
Link to repository removed for double-blind review.

\section*{AI Use}
\placeholder{} (counts toward the 20-page limit)

\bibliographystyle{ACM-Reference-Format}
\bibliography{paper}

\end{document}
